% !TEX root = ../main.tex
\lecture{1}{2026-09-14}{Tag 1}

\section{Admin stuff}
\begin{itemize}
    \item Online quiz 1x Woche Repetition (nicht Noten relevant) bis Montag 7.30 Uhr
    \item Definitionen und Sätze im Skript
    \item Alte Klausuren vefügbar zum lernen
    \item interaction points (nicht Noten relevant) für engagement und übungen
    \item 3 Challenges Mensa gutschein 
    \item Videos zum Thema der Woche immer 1 Woche früher vefügbar, machmal interessantes drin aber nicht prüfungsrelevant
    \item \textbf{Klausur relevant} Skript und Buch
    \subitem 30-40 Seiten pro Woche
    \subitem mit * Markiert bedeutet schwieriges thema nur für >5 Note relevant
    \subitem mit \# markiert ist nicht prüfungsrelevant
    \item Videos helfen vor der vorlesung um besser checken 
    \item Podcast nach vorlesung verfügbar
    \item \textbf{Musterlösungsweg} der Aufgaben der Woche am Ende der Woche verfügbar
    \subitem sonst Endresultate im Skript
    \item Tutorate sind kleinere übungsgruppen \textbf{heute ab 10 Uhr vefügbar}
    \item \textbf{Sprechstunden mit Profs} möglich oder mit Stundenen
    \item \textbf{KLAUSUR}
    \subitem keine Hilfsmittel? TS und Formelsammlung
    \item Voraussentzung WISSEN
    \subitem \textbf{Polynomdivision!} 
    \subitem \textbf{Ungleichungen!}
\end{itemize}

\section{Die Sprache der Mathematik}

\subsection{Logische Schlussfolgerungen}
\textbf{Aussage}

Ist ein Satz, der mit boolischer Werten überprüft werden kann (Ja oder Nein)

\begin{itemize}
    \item Implikation (A $\Rightarrow $ B)
    \subitem Es 12 Uhr ist $\Rightarrow$ die Glock läutet 
    \subitem Schlussfolgerung einer Aussage bei True oder False
    \item Umkehrung der Implikation (B $\Rightarrow$ A) ist NICHT immer wahr 
    \subitem die GLocke läutet $\Rightarrow$ es 12 Uhr ist 
    \item Kontraposition 
\end{itemize}

\textbf{Äquivalenz} (A $\Leftrightarrow B$)
Wenn A wahr dann B wahr, wenn B wahr dann A wahr

\subsection{Mengen}
Hat eine klare Definition einer Gruppe von Objekten, sodass bei gegebnem Objekt stets eindeutig feststellen zu welcher Menge es gehört.

\textbf{Definitionsmöglichkeiten}
\begin{itemize}
    \item Auflisten
    \subitem $ = {1,2,3...}$
    \item Eigenschaftbeschreibung
    \subitem $A = {a }$ 
\end{itemize}

\textbf{Mengenbeziehungen}
Teilmenge besteht wenn alle Elementen einer Menge in einer gleich ($\subseteq$) oder grösseren Menge ($\subset$ echte Teilmenge) enthalten ist.

\begin{itemize}
    \item A $\subset$ B
\end{itemize}

\textbf{Mächtigkeit einer Menge} (|A|)

Wenn Menge A endlich mit n-Elementen ist ihre Mächtikeit = n
\textbf{Aussage}
Wenn A = B dann immer |A| = |B|
\textbf{Umkehrung}
|A| = |B| nicht immer A = Buch

\textbf{Leeremengen}
Menge ohne Element aber ist TEilmenge aller anderen Mengen. Leer weil keine Elemente die Beschreibung erfüllen.

\textbf{Mengenoperationen}

\textbf{\textit{Kommutative Operationen}}

Vereinigung A $\cup$ B = B $\cup$ A

Schnitt A $\cap$ B = B $\cap$ A wenn A $\cap$ B = \{\}\} (leere Menge) \textit{Disjunkt}

\textbf{\textit{Nicht kommutative Ops}}
Differenz A \\ B oder, A ohne B, oder alle Elemente von A die NICHT in B enthalten sind

\textit{Komplement}
Wenn A $\subseteq B \Rightarrow A \setminus \textbf{Komplement} von B $ 


\subsubsection{Karthesisches Produkt} (A $\times$ B)

Neue menge aus \textbf{Tupel} (Paare) der Mengen A und B = {(-1,2), (-1,3),...} und = {(-1,3), (-1,2),...}